\begin{theorem}[Dini's theorem] \label{theorem:dini_theorem_for_monotone_pointwise_convergence_of_continuous_functions_on_a_compact_space_implies_uniform_convergence}
    \TODO{AI generated, verify}
    \TODO{Is there a notion of standard norm on a general commutative monoid or something that generalizes the notion of standard absolute value? And can this theorem be stated for more general spaces, not just totally ordered rings}
    Let $(X, \mathcal{T})$ be a \CrefAndHyperrefIfExist{definition:compact_topological_space}{compact topological space}, let $R$ be a commutative \CrefAndHyperrefIfExist{definition:ordered_ring}{totally ordered ring} equipped with its \CrefAndHyperrefIfExist{definition:standard_absolute_value_of_an_ordered_ring}{standard absolute value} $|\cdot|: R \to R_{\geq 0}$\CrefIfExists{definition:positive_cone_in_an_ordered_monoid}, and regard $R$ as an \CrefAndHyperrefIfExist{definition:extended_metric_on_a_set_valued_in_an_ordered_commutative_monoid_with_adjoined_infinity}{extended metric space} via the metric $d(a,b) = |a - b|$\CrefIfExists{definition:metric_induced_by_an_absolute_value_on_a_commutative_ring}.

    Let $(f_n)_{n=1}^\infty$ be a sequence of \CrefAndHyperrefIfExist{definition:continuous_map_of_topological_spaces}{continuous} functions $f_n : X \to R$ that converges \CrefAndHyperrefIfExist{definition:pointwise_convergence_of_a_sequence_of_functions_from_a_set_to_an_extended_metric_space_valued_in_ordered_commutative_monoid_with_adjoined_infinity}{pointwise} to a continuous function $f : X \to R$, and suppose further that the sequence $(f_n)_{n=1}^\infty$ is \CrefAndHyperrefIfExist{definition:monotone_sequence_of_elements_in_a_partially_ordered_set}{monotone} in the sense that for every $x \in X$, the sequence $(f_n(x))_{n=1}^\infty$ in $R$ is either monotonically increasing (i.e.\ $f_n(x) \leq f_{n+1}(x)$ for all $n$) or monotonically decreasing (i.e.\ $f_n(x) \geq f_{n+1}(x)$ for all $n$).

    Then the convergence $f_n \to f$ is \CrefAndHyperrefIfExist{definition:uniform_convergence_of_a_sequence_of_functions_from_a_set_to_an_extended_metric_space_valued_in_ordered_commutative_monoid_with_adjoined_infinity}{uniform} on $X$.
\end{theorem}

\begin{proof}
    \TODO{AI generated, verify}
    We treat the case where $(f_n)$ is monotonically decreasing, i.e.\ $f_n(x) \geq f_{n+1}(x)$ for all $n \in \bbN$ and $x \in X$; the increasing case follows by symmetry (consider $-f_n$ and $-f$, or equivalently apply the decreasing argument to the sequence $f - f_n$).

    We first define the auxiliary functions $h_n := f_n - f$. For each $n$, the function $h_n : X \to R$ is continuous because the map $(f_n, f) : X \to R \times R$ is continuous (for any open sets $U, V \subseteq R$, the preimage $(f_n, f)^{-1}(U \times V) = f_n^{-1}(U) \cap f^{-1}(V)$ is open in $X$, and products $U \times V$ generate the \CrefAndHyperrefIfExist{definition:product_topology_on_a_cartesian_product_of_topological_spaces}{product topology} on $R \times R$), the map $\sigma : R \times R \to R$ defined by $\sigma(a,b) = a - b$ is continuous by \Cref{proposition:commutative_ring_with_absolute_value_valued_in_ordered_semiring_has_continuous_algebraic_operations_under_metric_topology}, and the composition $h_n = \sigma \circ (f_n, f)$ is continuous by \Cref{proposition:composition_of_continuous_functions_on_topological_spaces_is_continuous}. By the monotonicity hypothesis and the pointwise convergence $f_n \to f$, we have $h_n(x) \geq h_{n+1}(x) \geq 0$ for all $n$ and $x$, and $h_n(x) \to 0$ pointwise.

    \TODO{prove that convergence of sequences preserves order}

    Fix $\varepsilon \in R_{>0}$. For each $n \in \bbN$, define the set
    $$
    U_n := \{x \in X : h_n(x) < \varepsilon\}.
    $$
    Since $h_n$ is continuous and takes non-negative values, the condition $h_n(x) < \varepsilon$ is equivalent to $|h_n(x)| < \varepsilon$, i.e.\ $h_n(x) \in B(0_R, \varepsilon)$ where $B(0_R, \varepsilon)$ is the open ball of radius $\varepsilon$ centered at $0_R$ in the metric topology on $R$\CrefIfExists{definition:topology_induced_by_an_extended_metric_valued_in_an_ordered_commutative_monoid_with_adjoined_infinity_and_open_ball_for_an_extended_metric}. Thus $U_n = h_n^{-1}(B(0_R, \varepsilon))$ is open in $X$.

    The sets $U_n$ are nested: because $h_n \geq h_{n+1}$, we have $U_n \subseteq U_{n+1}$. Indeed, if $h_n(x) < \varepsilon$, then $h_{n+1}(x) \leq h_n(x) < \varepsilon$, so $x \in U_{n+1}$.

    The family $\{U_n\}_{n=1}^\infty$ covers $X$. For any $x \in X$, the pointwise convergence $h_n(x) \to 0$ implies that there exists $N \in \bbN$ with $h_N(x) < \varepsilon$, i.e.\ $x \in U_N$.

    Since $X$ is compact and $\{U_n\}_{n=1}^\infty$ is an open cover, it admits a finite subcover. By the nesting property $U_1 \subseteq U_2 \subseteq \cdots$, any finite subcover is contained in $U_N$ for $N$ equal to the maximum of the indices appearing. Hence there exists $N \in \bbN$ such that $U_N = X$. Consequently, for every $x \in X$,
    $$
    |f_N(x) - f(x)| = h_N(x) < \varepsilon.
    $$

    For any $n \geq N$ and $x \in X$, the monotonicity $h_n(x) \leq h_N(x)$ implies $|f_n(x) - f(x)| = h_n(x) \leq h_N(x) < \varepsilon$. Thus $f_n \to f$ uniformly on $X$.
\end{proof}